
% ----------------------------------------------------------
% Anexos
% ----------------------------------------------------------
%
% ---
% Inicia os anexos
% ---
\begin{anexosenv}

% Imprime uma página indicando o início dos anexos
%\partanexos

% ---
\chapter{Calibração do tubo de Pitot}
\label{anexo:ensaio_pitot}
% ---

Ensaio realizado pela Equipe Albatroz no Instituto de Pesquisas Tecnológicas
(IPT), em túnel de vento, para o tubo de Pitot instrumentado com o sensor de
pressão diferencial MPX5002DP. A leitura foi feita por conversor A/D de
10~bits com referência de \SI{5}{\volt} (\SI{4,89}{\milli\volt} por LSB), em
seis patamares de velocidade, com densidade do ar de
\SI{1,110}{\kilo\gram\per\meter\cubed}.

\begin{table}[H]
    \centering
    \caption{Leituras por patamar de velocidade}
    \label{tab:pitot_dados_brutos}
    \small
    \begin{tabular}{ccccc}
        \toprule
        \textbf{$v_{ref}$} & \textbf{Amostras} & \textbf{ADC médio} &
        \textbf{Desvio} & \textbf{Tensão média} \\
        {[\si{\meter\per\second}]} & & [LSB] & [LSB] & [\si{\volt}] \\
        \midrule
        1,00  & 55  & 546,02 & 0,216 & 2,6687 \\
        2,89  & 155 & 545,38 & 1,959 & 2,6657 \\
        5,60  & 159 & 547,20 & 0,886 & 2,6745 \\
        8,40  & 149 & 549,86 & 0,775 & 2,6875 \\
        11,57 & 155 & 554,51 & 0,868 & 2,7103 \\
        14,80 & 186 & 561,07 & 2,445 & 2,7424 \\
        \bottomrule
    \end{tabular}

    \vspace{2pt}
    {\small Fonte: Equipe Albatroz (2026)}
\end{table}

O ajuste do modelo quadrático $ADC = k\,v^{2} + ADC_0$ fornece
$k = \num{0,071521}$~LSB/(\si{\meter\per\second})$^2$ e
$ADC_0 = \num{545,14}$~LSB, com $R^2 = \num{0,9944}$. Separando as constantes
do sensor das do conversor, obtêm-se a sensibilidade
$S = \SI{629,8}{\micro\volt\per\pascal}$ e a saída com escoamento nulo
$V_0 = \SI{2,66442}{\volt}$, que são as grandezas transportadas para o
\textit{firmware} da SIBA (\autoref{sec:calibracao_pitot}).

Acima de \SI{5}{\meter\per\second} o erro do ajuste fica abaixo de
\SI{4,2}{\percent}. Nos dois primeiros patamares a leitura varia menos de
\SI{1}{LSB} e o sinal é dominado pela quantização do conversor de 10~bits.

\end{anexosenv}

